\documentclass[11pt]{article}
\usepackage[margin=1in]{geometry}
\usepackage[T1]{fontenc}
\usepackage[utf8]{inputenc}
\usepackage{lmodern}
\usepackage{setspace}
\usepackage{parskip}
\usepackage{amsmath}
\usepackage{amssymb}
\usepackage{booktabs}
\usepackage{hyperref}
\usepackage{enumitem}

\hypersetup{
  colorlinks=true,
  linkcolor=black,
  urlcolor=blue,
  citecolor=black,
  pdftitle={MDPinn Training and Evaluation Refactor Plan},
  pdfauthor={OpenAI Codex}
}

\setstretch{1.1}
\setlist[itemize]{leftmargin=1.5em}

\begin{document}

\title{MDPinn Training/Evaluation Refactor and Rollout-Aware Training Plan}
\date{May 1, 2026}
\maketitle

\section{Purpose}

This document replaces the earlier manuscript-style outline with an engineering note that is directly aligned with the current state of the repository. The immediate objective is to make the workflow clean and reproducible: training should produce checkpoints, evaluation should consume checkpoints, and the graph-generation and model-selection logic should be rerunnable without retraining. The second objective is to document why the current training setup does not line up with the rollout behavior we actually care about, and to specify the next training changes required to fix that mismatch.

\section{Current Repository State}

The repository currently contains a strong training core and several useful evaluation utilities, but the workflow has been blurred by convenience scripts that mix model production with post-processing. In particular:

\begin{itemize}
  \item \texttt{src/run\_training.py} has been used as a high-level entry point for training.
  \item \texttt{src/compare\_models.py}, \texttt{src/compare\_rollouts.py}, and \texttt{src/rollout\_nve.py} already provide most of the evaluation logic needed for static metrics, rollout summaries, and drift plots.
  \item \texttt{src/experiment\_suite.py} can train, compare, and run rollout post-processing in one pass, which is useful for large sweeps but obscures the conceptual boundary between training and evaluation.
\end{itemize}

That ambiguity matters because the project is now at the point where checkpoints must be inspected repeatedly under revised diagnostics. If evaluation is coupled to training, then every new graph or metric change becomes more expensive than necessary and makes it harder to reason about which result came from which checkpoint.

\section{Completed Workflow Refactor}

The repository has now been reorganized around a clean two-phase workflow.

\subsection{Training Phase}

Training is handled by \texttt{src/run\_training.py}. This runner is now intentionally training-only. It supports:

\begin{itemize}
  \item standard training,
  \item physics-informed training,
  \item Optuna tuning.
\end{itemize}

It no longer presents ``train and compare'' as a primary mode. The purpose of this change is to make checkpoint production independent from checkpoint analysis.

\subsection{Evaluation Phase}

Post-training analysis is handled by a new dedicated entry point: \texttt{src/run\_evaluation.py}. This script supports two use cases:

\begin{itemize}
  \item evaluating a single checkpoint to produce held-out frame metrics, single-model plots, and rollout diagnostics;
  \item evaluating a pair of checkpoints to produce framewise comparison plots and rollout-comparison artifacts.
\end{itemize}

This separation is necessary because evaluation is now expected to evolve faster than training. We need to be able to revisit existing checkpoints with better rollout-aware metrics, different graph settings, and revised model-selection logic without paying the cost of retraining.

\subsection{Experiment Configuration Defaults}

The example experiment configuration has also been updated so that post-training comparison and rollout generation are no longer enabled by default. The automation hooks remain available for explicit use, but evaluation is now treated as opt-in post-processing rather than part of the default training path.

\section{Why the Current Training Objective Does Not Match the Rollout Goal}

The central technical problem in the current setup is that the model is primarily trained as a supervised framewise regressor, while we are evaluating it as a closed-loop molecular dynamics surrogate.

\subsection{What the Model Is Optimizing Right Now}

At present, the dominant training objective is supervised fitting of:

\begin{itemize}
  \item energies,
  \item forces.
\end{itemize}

The force term is heavily weighted relative to the energy term. In the present configuration, the model is mainly being asked to achieve good local agreement with reference labels on dataset frames. In delta-learning mode, it is learning residual energy and residual force targets relative to the analytic baseline rather than the full observables directly, but this does not change the basic character of the optimization problem: the primary target remains local supervised accuracy on data-manifold configurations.

\subsection{What the Rollout Evaluation Asks the Model to Do}

The rollout code uses velocity-Verlet integration and repeatedly applies the model's own forces to advance the system in time. This is a qualitatively different regime. Good rollout behavior requires not only low local force error, but also:

\begin{itemize}
  \item a smooth and conservative effective energy surface between labeled configurations,
  \item robustness to repeated self-generated integration error,
  \item acceptable behavior after the trajectory drifts away from the exact training manifold.
\end{itemize}

These are not guaranteed by strong one-step supervised performance. A model can interpolate labeled frames reasonably well while still producing trajectories with large long-horizon drift.

\subsection{The Core Mismatch}

The current system is therefore solving:
\[
\text{fit snapshot labels well}
\]
while we actually care about:
\[
\text{remain stable under repeated model-driven integration.}
\]

This mismatch explains why rollouts can still look poor even when static metrics appear respectable.

\section{Why the Existing Physics Terms Have Not Mattered Enough}

The repository already contains momentum and NVE-related regularization, but in practice those terms have not yet become decisive in training.

\subsection{They Are Soft Penalties, Not Primary Objectives}

The current physics-informed model adds momentum and NVE losses on top of the supervised objective. The structure is:
\[
\mathcal{L}_{\text{total}}
=
\mathcal{L}_{\text{sup}}
+
\lambda_{\text{mom}} \mathcal{L}_{\text{mom}}
+
\lambda_{\text{NVE}} \mathcal{L}_{\text{NVE}}
+
\lambda_{\text{PBC}} \mathcal{L}_{\text{PBC}}.
\]

That means the training loop still treats snapshot regression as the main job and physics as regularization. This is not wrong, but it means the auxiliary terms must have meaningful numerical scale if they are going to influence the optimizer.

\subsection{The Raw NVE Signal Is Too Small}

The observed training logs show that the NVE loss is numerically tiny compared with the supervised loss. If the raw NVE magnitude is effectively near zero, increasing its weight has limited effect. A near-zero term multiplied by a moderate scalar is still near zero.

\subsection{The NVE Term Is Sparse}

The current NVE penalty is not applied on every batch. It is gated by:

\begin{itemize}
  \item an application frequency,
  \item a warmup period,
  \item a ramp schedule.
\end{itemize}

That scheduling was sensible during development, but it further reduces the influence of a term that is already small.

\subsection{The NVE Formulation Is Indirect}

The present NVE-style regularizer is evaluated on short trajectory windows built from reference data rather than on short learned rollouts generated by the model itself. This makes it a useful on-manifold smoothness prior, but it does not directly train the precise failure mode revealed by evaluation, namely long-horizon closed-loop instability under the deployed integrator.

\section{Why NVE Is a Difficult Objective}

The NVE objective is valuable, but it is structurally harder than ordinary framewise supervision.

\begin{itemize}
  \item It is trajectory-level rather than pointwise.
  \item Small local force biases can compound into large long-horizon drift.
  \item The physically relevant quantity is total-energy stability under integration, which depends on both the learned surface and the numerical rollout procedure.
  \item The signal is more expensive to compute because it requires multiple force evaluations across a trajectory segment rather than one forward pass on a snapshot.
\end{itemize}

This is why rollout-aware training has to be designed deliberately. It cannot be treated as a minor add-on and expected to dominate a snapshot-driven objective automatically.

\section{Planned Training Changes}

The next training changes should be implemented in stages, because they affect both optimization and model-selection logic.

\subsection{1. Introduce Rollout-Aware Validation}

The first change is to define a validation metric that actually reflects the deployment goal. The proposed primary validation score is:
\[
\text{median over rollouts of } \frac{1}{T} \sum_t |E(t)-E(0)|.
\]

This should be paired with:

\begin{itemize}
  \item median max absolute drift per rollout,
  \item rollout failure rate.
\end{itemize}

A practical scalar checkpoint score is:
\[
\text{score}
=
\text{median mean } |dE|
+
0.5 \times \text{median max } |dE|
+
50 \times \text{failure rate}.
\]

This change is necessary because model selection currently prefers low validation snapshot loss, not good rollout behavior.

\subsection{2. Select Checkpoints Using Rollout Metrics}

Once rollout-aware validation is available, checkpointing and early stopping must monitor that score instead of \texttt{val\_total\_mse\_loss}. Snapshot metrics should still be logged for diagnostics, but they should no longer be the sole decision rule for choosing the ``best'' model.

\subsection{3. Replace Reference-Trajectory NVE with Learned-Rollout NVE}

The NVE training term should be aligned with deployment. Instead of only evaluating drift along short windows of reference positions, the training loop should:

\begin{itemize}
  \item initialize from real training frames,
  \item reconstruct initial velocities from adjacent frames,
  \item integrate the model forward for a short learned rollout,
  \item penalize total-energy drift on that learned rollout.
\end{itemize}

This is necessary because the current NVE term is too indirect. The rollout failures we care about occur under model-generated trajectories, so the training signal should directly reflect that regime.

\subsection{4. Fix the Scale of the NVE Term}

Before tuning weights aggressively, the raw NVE term must be made numerically meaningful. Possible changes include:

\begin{itemize}
  \item using absolute drift in eV instead of a normalization that crushes the signal,
  \item using a controlled normalization scheme whose scale remains visible in training,
  \item logging raw and weighted physics-term magnitudes every epoch.
\end{itemize}

This is necessary because without proper scale, weight tuning becomes mostly cosmetic.

\subsection{5. Increase the Effective Importance of Physics Terms}

After the formulation and scale are fixed, the following knobs should be revisited:

\begin{itemize}
  \item lower \texttt{nve\_freq},
  \item reduce or remove long warmup,
  \item shorten rollout length if needed to afford more frequent application,
  \item retune \texttt{nve\_weight} and \texttt{momentum\_weight}.
\end{itemize}

This is necessary because even a better loss will not matter if it is applied too rarely or too late.

\subsection{6. Keep Static Metrics as Secondary Diagnostics}

Energy and force errors should continue to be reported. They remain important sanity checks. However, they should be interpreted as local fit diagnostics rather than as the final authority on whether a checkpoint is suitable for molecular dynamics rollout.

\section{Implementation Order}

The recommended order of work is:

\begin{enumerate}
  \item keep training and evaluation separate;
  \item add rollout-aware validation utilities;
  \item switch checkpoint selection to rollout-aware validation;
  \item implement learned-rollout NVE training;
  \item recalibrate NVE scale;
  \item retune physics frequency, schedule, and weights;
  \item compare old and new training setups with a small controlled sweep.
\end{enumerate}

This order is necessary because changing weights before fixing the validation target and NVE formulation would make the optimization harder to interpret and easier to misdiagnose.

\section{Updated Workflow Instructions}

\subsection{Training}

Use \texttt{src/run\_training.py} only to create checkpoints.

\begin{verbatim}
python src/run_training.py --mode standard --molecule aspirin
python src/run_training.py --mode physics --molecule aspirin
\end{verbatim}

\subsection{Single-Checkpoint Evaluation}

Use \texttt{src/run\_evaluation.py} with \texttt{--checkpoint} to generate held-out metrics, single-model plots, and rollout diagnostics for one checkpoint.

\begin{verbatim}
python src/run_evaluation.py ^
  --checkpoint checkpoints/physics_informed/best_model.ckpt ^
  --molecule aspirin ^
  --output-dir results/eval/physics_best
\end{verbatim}

\subsection{Paired Evaluation}

Use \texttt{src/run\_evaluation.py} with \texttt{--standard} and \texttt{--candidate} to generate framewise comparison plots and rollout-comparison summaries.

\begin{verbatim}
python src/run_evaluation.py ^
  --standard checkpoints/standard/best_model.ckpt ^
  --candidate checkpoints/physics_informed/best_model.ckpt ^
  --molecule aspirin ^
  --output-dir results/eval/std_vs_phys ^
  --standard-label Standard ^
  --candidate-label Physics
\end{verbatim}

\section{Conclusion}

The repository now has the infrastructure needed for a cleaner research loop: train once, evaluate repeatedly, and revise the diagnostics without entangling them with model production. That separation is a prerequisite for the next phase of work, because the main remaining issue is not merely code organization. It is a methodological mismatch between snapshot-oriented training and rollout-oriented use. The planned rollout-aware validation and learned-rollout NVE changes are necessary precisely because the current objective does not place enough optimization pressure on the behavior we actually care about.

\end{document}
